\chapter{Общий калибровочный якорь спектрального момента}
\label{ch:v7-common-gauge-f0-anchor-gate}

\section{Условный словарь коэффициентов}

Предыдущий гейт получил
\begin{equation}
 C_0=\frac{f_0}{8\pi^2},
 \qquad
 \lambda_E=\frac1{8C_0}=\frac{\pi^2}{f_0}.
 \label{eq:v7-gauge-anchor-edge-quartic}
\end{equation}
Пусть физический калибровочный генератор имеет полный конечный индекс
\begin{equation}
 q_G=\Tr_{\mathcal H_F}Q^2.
 \label{eq:v7-gauge-anchor-total-index}
\end{equation}
В конвенции heat-kernel словаря Тома IV коэффициент кривизны равен
\begin{equation}
 \frac{C_0q_G}{3}F_{\mu\nu}F^{\mu\nu}.
 \label{eq:v7-gauge-anchor-curvature-coefficient}
\end{equation}
Сопоставление с $F^2/(4g^2)$ условно даёт
\begin{equation}
 f_0=\frac{6\pi^2}{q_Gg^2},
 \qquad
 \boxed{\lambda_E=\frac{q_G}{6}g^2.}
 \label{eq:v7-gauge-anchor-conditional-relation}
\end{equation}
Для старого относительного $U(1)$-словаря с $q_G=2$ это стало бы
\begin{equation}
 \lambda_E=\frac13g^2.
 \label{eq:v7-gauge-anchor-old-u1-relation}
\end{equation}
Таким образом, алгебраическое сокращение $f_0$ не является проблемой.
Проверяемый вопрос состоит в том, определён ли тот же $q_G$ текущим
рёберным носителем.

\section{Два неравных конечных следа}

Hodge-родитель использует редуцированное пространство меток
\begin{equation}
 \mathcal E_{\rm new}^{\rm red}\simeq\mathbb C^{11},
 \qquad
 \Tr_{\rm red}|e\rangle\langle e|=1
 \quad\text{для каждого ребра }e.
 \label{eq:v7-gauge-anchor-reduced-edge-trace}
\end{equation}
Именно этот незавешенный след дал одинаковые квартичные коэффициенты и
точный селектор шести рёбер.

Физическая калибровочная кривизна действует иначе. На блоке
$\operatorname{Hom}(R_s,R_t)$ индуцированный абелев заряд равен
\begin{equation}
 q_e=Y_t-Y_s,
 \qquad
 \Tr_e F^2\ \mathrel{\propto}\
 \dim(R_e)q_e^2F^2,
 \label{eq:v7-gauge-anchor-induced-charge}
\end{equation}
а для неабелевых факторов вместо $q_e^2$ входят индексы представлений.
Следовательно, физический след зависит от раскрытых калибровочных
кратностей, а не только от числа рёбер \cite{v7-gauge-anchor-beyond-sm}.

Уже на шести выбранных блоках абелевы разности равны
\begin{equation}
 \begin{array}{c|cccccc}
 e&L_LY_R&Q_LY_R&X_LX_R&X_Le_R&X_Lu_R&Y_LY_R\\ \hline
 Y_t-Y_s&0&-\frac23&0&0&\frac53&0
 \end{array}.
 \label{eq:v7-gauge-anchor-selected-charges}
\end{equation}
Четыре канала являются синглетными по этому $U(1)$, тогда как два имеют
разные ненулевые веса. Даже до цветовых и слабых множителей
\begin{equation}
 \sum_{e\in E_*}(Y_t-Y_s)^2=\frac{29}{9},
 \qquad
 \{q_e^2:e\in E_*\}=\left\{0,\frac49,\frac{25}{9}\right\}.
 \label{eq:v7-gauge-anchor-nonuniform-index}
\end{equation}
Поэтому незавешенный Hodge-след не является физическим gauge-следом.

\section{Почему прямое присоединение не закрывает вопрос}

Можно формально взять блочную сумму физического $H_{15}$ и
вспомогательного комплекса стрелок и применить одну функцию отсечения.
Однако тогда калибровочный индекс равен
\begin{equation}
 q_G^{\rm tot}=q_G^{H_{15}}+q_G^{\rm edge},
 \label{eq:v7-gauge-anchor-total-block-index}
\end{equation}
где $q_G^{\rm edge}$ требует полного представления каждого
калибровочно-ковариантного блока. Редуцированная модель $\mathbb C^{11}$
эти данные стёрла. Положить $q_G^{\rm edge}=0$ нельзя: смешанные рёбра
явно заряжены. Использовать только старое $q_G^{H_{15}}=2$ также нельзя,
поскольку тогда Hodge-блок присутствует в потенциале, но удалён из
кривизностного следа.

Кроме того, нечётный оператор рёбер бимодульно линеен и имеет
$\Omega^1_{D_E}(\mathcal A_F)=0$. Поэтому он был допущен как
ассоциированный полевой комплекс, но не как обычная физическая конечная
спектральная тройка. Теория произведений допускает индуцированные связности
на таких пучках, но не отождествляет их след с фермионным следом без явного
модуля и представления \cite{v7-gauge-anchor-product}.

\begin{theorem}[Условное сокращение и запрет старого якоря]
При заранее фиксированном полном калибровочном индексе $q_G$ один
спектральный момент действительно даёт
$\lambda_E=q_Gg^2/6$. Но текущий редуцированный Hodge-носитель считает
одинаково одиннадцать меток рёбер и не определяет физические индексы их
представлений. Поэтому значение $q_G=2$ из прежнего физического словаря
нельзя перенести на новый product-оператор; отношение
$\lambda_E=g^2/3$ остаётся условным.
\end{theorem}

\begin{nogocriterion}[Запрет смешения следов]
Нельзя нормировать gauge-кинетику полным физическим следом, а Hodge-потенциал
--- незавешенным следом по меткам рёбер и называть их одним спектральным
действием. Допустим либо полный калибровочно-взвешенный рёберный носитель с
повторным вычислением потенциала и гессиана, либо явное признание двух
разных функционалов.
\end{nogocriterion}

\begin{protocol}
\begin{enumerate}
 \item формальное сокращение $f_0$: \textbf{получено};
 \item условное отношение: \textbf{$\lambda_E=q_Gg^2/6$};
 \item старый словарь $q_G=2$: \textbf{$\lambda_E=g^2/3$ условно};
 \item равные Hodge-веса рёбер: \textbf{да};
 \item равные gauge-индексы рёбер: \textbf{нет};
 \item синглетных выбранных каналов по $U(1)$: \textbf{$4$};
 \item заряженных выбранных каналов по $U(1)$: \textbf{$2$};
 \item полный $q_G^{\rm edge}$: \textbf{редуцированным носителем не задан};
 \item перенос прежнего $q_G=2$: \textbf{запрещён};
 \item итог: \textbf{условная формула и no-go прямого gauge-якоря};
 \item следующий гейт: \textbf{полный gauge-взвешенный носитель рёбер}.
\end{enumerate}
\end{protocol}

Машинный сертификат сохранён в
\path{s2t_v7_common_gauge_f0_anchor_gate_results.json}.

\begin{thebibliography}{9}
\bibitem{v7-gauge-anchor-spectral-action}
A.~H.~Chamseddine, A.~Connes,
\emph{The Spectral Action Principle},
Communications in Mathematical Physics 186 (1997), 731--750;
arXiv:hep-th/9606001.
\bibitem{v7-gauge-anchor-beyond-sm}
W.~D.~van Suijlekom,
\emph{Going beyond the Standard Model with noncommutative geometry},
arXiv:1211.0825.
\bibitem{v7-gauge-anchor-product}
S.~Brain, B.~Mesland, W.~D.~van Suijlekom,
\emph{Gauge Theory for Spectral Triples and the Unbounded Kasparov Product},
Journal of Noncommutative Geometry 10 (2016), 135--206;
arXiv:1306.1951.
\end{thebibliography}